\section{心理尺度を作る}
\subsection{授業内容}
\subsubsection{科目の中でのこのコマの位置づけ}

\subsubsection{コマ主題細目}
\begin{description}
    \item[サーストンの等現間隔法]
    \item[リッカーとのシグマ法]
    \item[尺度を評価する] IT相関，内的整合性信頼性

\end{description}
\subsubsection{キーワード}
\begin{itemize}
\item サーストンの等現間隔法
\item リッカートのシグマ法
\item 信頼性
\item 妥当性
\item ワーディングの問題
\end{itemize}
\subsection{授業情報}
\paragraph{コマの展開方法}
講義
\subsubsection{予習・復習課題}
\paragraph{予習}
\paragraph{復習}
